\section{MI-06: no finite two-unitary constant for the arithmetic symmetric modulus}
\label{sec:mi06}

Define the arithmetic symmetric modulus by
\[
 \symmod(X)=\frac{|X|+|X^*|}{2}.
\]
The target proposes that for every $X,Y$ there are unitaries $U,V$ with
\[
 \symmod(X+Y)\leq\sqrt2\bigl(U\symmod(X)U^*+V\symmod(Y)V^*\bigr).
\]
The following result disproves that assertion and the more general finite-constant question in \cite[Conjecture 1.6 and equation (1.4)]{Zhang2026}.

\begin{theorem}\label{thm:mi06}
Even in dimension three, there is no finite nonnegative constant $C$ such that every pair $X,Y\in\M_3(\C)$ admits unitaries $U,V$ satisfying
\begin{equation}\label{eq:mi06-target}
 \symmod(X+Y)\leq C\bigl(U\symmod(X)U^*+V\symmod(Y)V^*\bigr).
\end{equation}
The failure is witnessed by real rank-one matrices.
\end{theorem}

\begin{proof}
For $t>0$, let
\[
 X=e_1(e_1+te_2)^*
 =\begin{pmatrix}1&t&0\\0&0&0\\0&0&0\end{pmatrix},\qquad
 Y=-(e_1+te_3)e_1^*
 =\begin{pmatrix}-1&0&0\\0&0&0\\-t&0&0\end{pmatrix}.
\]
Put $s=\sqrt{1+t^2}$. If $Z=uv^*$ has rank one, then
\[
 |Z|=\frac{\|u\|}{\|v\|}vv^*,\qquad
 |Z^*|=\frac{\|v\|}{\|u\|}uu^*.
\]
Applying this identity shows that both $\symmod(X)$ and $\symmod(Y)$ have eigenvalues
\[
 a=\frac{s+1}{2},\qquad b=\frac{s-1}{2},\qquad0.
\]
On the other hand,
\[
 X+Y=t(e_1e_2^*-e_3e_1^*)
\]
has right and left moduli $t\diag(1,1,0)$ and $t\diag(1,0,1)$, respectively. Hence
\begin{equation}\label{eq:mi06-sum}
 \symmod(X+Y)=\diag(t,t/2,t/2)\geq(t/2)I_3.
\end{equation}

Fix arbitrary unitaries $U,V$. Choose a unit vector $w$ orthogonal to a top eigenvector of $U\symmod(X)U^*$ and to a top eigenvector of $V\symmod(Y)V^*$. Such a vector exists because the ambient dimension is three. The eigenvalue lists above imply
\[
 w^*\bigl(U\symmod(X)U^*+V\symmod(Y)V^*\bigr)w\leq 2b=s-1.
\]
Combining this with \eqref{eq:mi06-sum}, any domination \eqref{eq:mi06-target} would force
\[
 \frac t2\leq C(s-1),\qquad
 C\geq\frac{t}{2(\sqrt{1+t^2}-1)}
 =\frac{\sqrt{1+t^2}+1}{2t}.
\]
The final expression tends to $+\infty$ as $t\downarrow0$.
\end{proof}

\begin{corollary}[A small exact counterexample to $\sqrt2$]\label{cor:mi06-rational}
Take $t=3/4$ in the preceding construction. Then
\[
 \spec\symmod(X)=\spec\symmod(Y)=\{9/8,1/8,0\},
 \qquad \symmod(X+Y)=\diag(3/4,3/8,3/8).
\]
For arbitrary $U,V$, the quadratic form of the proposed right-hand side on the vector $w$ used above is at most $\sqrt2/4<3/8$. Thus the conjectured domination fails.
\end{corollary}

For direct verification, the two individual symmetric moduli in this example are rational matrices:
\[
 \symmod(X)=\begin{pmatrix}41/40&3/10&0\\3/10&9/40&0\\0&0&0\end{pmatrix},\qquad
 \symmod(Y)=\begin{pmatrix}41/40&0&3/10\\0&0&0\\3/10&0&9/40\end{pmatrix}.
\]
This obstruction concerns order domination by two separately rotated summands. It does not contradict a unitarily invariant norm inequality for their sum, and it does not identify the arithmetic modulus with the quadratic symmetric modulus.
